\begin{papercard}{aoir2019}
\cf{Stages} \stagetag{D} (collecting store pages, scraped listings and screens), \stagetag{W} (walkthrough as research on a platform), \stagetag{L} (use of machine-learning and LLM tools in analysis), \stagetag{R} (documentation, sharing and retention), \stagetag{S} (ethics of the design). franzke, Bechmann, Zimmer, Ess and the AoIR Ethics Working Group, ``Internet Research: Ethical Guidelines 3.0'' (Association of Internet Researchers, approved unanimously by the membership on 6 October 2019); 83 PDF pages \pc{aoir2019}{1}{Unanimously approved by the AoIR membership October 6, 2019}. The main text (PDF pages 3--25) is a question-led framework; companion resources follow: 6.1 AI and machine learning (pp.\ 34--50), 6.2 academy--industry partnerships and corporate data, 6.3 feminist research ethics, 6.4 an ``impact model'', and appendices including operational security. This card summarises the main text and 6.1; 6.2--6.4 and the appendices were not read in detail. Page numbers below are PDF pages (the document's own page numbers are one lower).

\cf{Orientation} IRE 3.0 builds on IRE 1.0 (2002) and 2.0 (2012). It is ``premised on primary commitments to ethical pluralism and cross-cultural awareness, coupled with the experientially-grounded view that ethics starts by asking and answering critical questions'' rather than applying rules \pc{aoir2019}{3}{IRE 3.0 is premised on primary commitments to ethical pluralism and cross-cultural awareness, coupled with the experientially-grounded view that ethics starts by asking and answering critical questions}. For review boards the authors argue that ongoing ethical reflection ``might be more helpful and beneficial in the long term for society than now restricting research'' \pc{aoir2019}{3}{An ongoing ethical reflection might be more helpful and beneficial in the long term for society than now restricting research}. The primary norms come from the Belmont Report: respect for persons, beneficence, justice \pc{aoir2019}{5}{as initially rooted in The Belmont Report}.

\cf{Process approach} Guidelines develop ``from the bottom up in a case-by-case based approach without making a priori judgements whether some research per se is unethical'' \pc{aoir2019}{5}{one that aims to develop guidelines from the bottom up in a case-by-case based approach without making a priori judgements whether some research per se is unethical}. ``Ethics is method -- method is ethics'': method choices raise ethical issues and revised methods require revisiting ethical assumptions \pc{aoir2019}{5}{ethics is method method is ethics}, which counters seeing ethics as a ``one-off tick-box exercise'' \pc{aoir2019}{5}{a common presumption of ethics as something of a one-off tick-box exercise that is primarily an obstacle to research}. Frameworks differ: continental and Scandinavian approaches ``are strongly deontological'', US/UK ones lean utilitarian \pc{aoir2019}{6}{European and Scandinavian approaches to research ethics in general and privacy matters in particular are strongly deontological}; ethical pluralism connects local practice to shared norms. The result is ``guidelines, not recipes'' \pc{aoir2019}{7}{the best we can do is develop guidelines, not recipes}.

\cf{Guiding questions from IRE 1.0 and 2.0} Where does the interaction under study take place and what ethical expectations does the venue set \pc{aoir2019}{8}{Where does the inter/action, communication, etc. under study take place}; who are the subjects; what significant risks and benefits does the research entail \pc{aoir2019}{8}{What ethically significant risks does the research entail for the subject(s)}. Users often expected exchanges in quasi-public venues to be private \pc{aoir2019}{8}{But they often nonetheless expected that these exchanges were somehow private}. Later questions concern data management and storage, breach and re-identification risk, and the danger that future tools (automated text analysis, face recognition) undo anonymisation.

\cf{Stages of research and consent} Stages: initial design, data acquisition, analysis, dissemination, close of project \pc{aoir2019}{10}{IRE 3.0 extends this approach by way of a more extensive taxonomy of stages}. Analysis can re-identify data through aggregation of datasets \pc{aoir2019}{10}{Analyses, including assessment of how use of particular techniques, formulas, or instruments may re-identify data through aggregation of multiple data sets}. Consent is ``manifestly impracticable'' in big-data projects \pc{aoir2019}{11}{Such consent is manifestly impracticable in the case of Big Data projects, however, resulting in a serious ethical dilemma}; mitigations include pseudonymisation with keys kept apart \pc{aoir2019}{11}{Most commonly, researchers pseudonymize their data separating keys from the actual dataset}, and asking consent only at dissemination for quotes that could identify people. A footnote asks whether publishing an exact quote is necessary \pc{aoir2019}{12}{why would publication of an exact quote be necessary}. Data minimisation: collect ``no more data, especially sensitive and personal information \ldots\ than is strictly necessary'' \pc{aoir2019}{13}{no more data, especially sensitive and personal information, is collected than is strictly necessary}.

\cf{Legal and platform issues} NESH guidelines are praised for combining individual and relational views and for foregrounding ``Nissenbaum's account of privacy as contextual integrity'' \pc{aoir2019}{14}{foreground Nissenbaum’s account of privacy as contextual integrity}. Researchers ``will increasingly confront terms of use that dictate the conditions under which such research activities can take place, if at all'' \pc{aoir2019}{15}{researchers will increasingly confront terms of use that dictate the conditions under which such research activities can take place, if at all}; Facebook and LinkedIn prohibit automated scraping \pc{aoir2019}{15}{Facebook and LinkedIn prohibit the automated scraping of data from their platform}. In the US, ignoring terms may touch the Computer Fraud and Abuse Act, and whether terms bind academic researchers ``remains debatable'' \pc{aoir2019}{16}{In the United States, failure to follow such terms may violate the Computer Fraud and Abuse Act}. Fake accounts are a form of deceptive research \pc{aoir2019}{16}{developing fake profiles is a form of deceptive research}. GDPR is named as the main legal frame in Europe.

\cf{Subjects and harm} ``A primary ethical imperative is to avoid harm'' to subjects and researchers \pc{aoir2019}{18}{A primary ethical imperative is to avoid harm - to subjects as well to researchers}; the greater the vulnerability, the greater the obligation \pc{aoir2019}{18}{the greater the vulnerability of our subjects, the greater our responsibility and obligation to protect them from likely harms}. Listed considerations: downstream harms \pc{aoir2019}{18}{Downstream harms or harms after the fact}, minors, politically sensitive research, women, people with disabilities, and special emotional states such as grief, trauma and illness. Accountability, trust and transparency are named as key concepts \pc{aoir2019}{19}{accountability, trust, and transparency as also critical to good research practice and ethics}.

\cf{Data handling questions} Questions on managing and storing data and on consequences of stripping identifiers \pc{aoir2019}{20}{What are the potential ethical consequences of stripping data of personally identifiable information}; ``the best current practices and techniques can only de-identify data \ldots\ not perfectly anonymize'' \pc{aoir2019}{20}{best current practices and techniques can only de-identify data, i.e., not perfectly anonymize data}. Data types differ (aggregated, derived, purchased) \pc{aoir2019}{21}{Differences between different types of data (aggregated sets, derived sets, purchased sets}; data minimisation conflicts with inductive big-data aims; sharing risks are illustrated by the OKCupid release \pc{aoir2019}{21}{OKCupid data release case where supposedly anonymized profiles were easily re-linked to persons}.

\cf{Algorithms and learning models (main text)} ``How can we make sure that our models and subsequent results are adequately documented in order to replicate our results'' \pc{aoir2019}{22}{How can we make sure that our models and subsequent results are adequately documented in order to replicate our results}; the layers of cleaning, classification and re-training ``require new standards for documentation'' \pc{aoir2019}{22}{The many layers and iterations of cleaning/pre-processing, classification, training/retraining in order to predict a trend, tendency or pattern require new standards for documentation}. Reverse engineering of closed algorithms faces scraping limits and IP protection \pc{aoir2019}{22}{Other strands of research within internet research use reverse engineering in order to understand closed algorithms}. Crowdsourced labelling raises pay and ``labelling quality'' concerns \pc{aoir2019}{23}{there are also ethical concerns regarding the labelling quality}. The closing advice: ``talk things over with colleagues and friends'' \pc{aoir2019}{24}{one of the most important ethical techniques to be recommended is one of the simplest: talk things over with colleagues and friends}.

\cf{Companion 6.1: AI and machine learning} Organised by phase. \emph{Design}: ethical dimensions conflict (more data to reduce bias versus privacy), so make choices explicit; state aims, since ``emerging risks will be judged against these in a balancing exercise'' \pc{aoir2019}{38}{Emerging risks will be judged against these in a balancing exercise}; combine an ex ante hypothesis with documentation of choices \pc{aoir2019}{39}{We suggest that researchers who use AI systems in their research consider a hybrid of an ex ante hypothesis (still making room for exploration) as well as documenting actions taken and choices made throughout the process}. Encryption improves security, not privacy \pc{aoir2019}{40}{applying technological solutions such as encryption are often mistakenly classed as efforts to improve privacy, while they instead provide more security}; ask whether data points are necessary and not excessive \pc{aoir2019}{40}{Are the identified data points necessary, relevant and not excessive in relation to the research aim}. \emph{Analysis}: training-data provenance and bias; the time and assumptions spent on cleaning should be reported \pc{aoir2019}{42}{The time spent on cleaning a dataset and the assumptions that go into this process should be communicated more clearly in the resulting research paper}; on automation, ``is there a strong incentive for the researcher to take the automated inferences as a base truth?'' \pc{aoir2019}{43}{Is there a strong incentive for the researcher to take the automated inferences as a base truth}; can the researcher override automated inferences and document that \pc{aoir2019}{43}{Can the researcher override the automated inferences procedure, how will this be documented and justified for later reproducibility}; account for false positives and negatives \pc{aoir2019}{43}{To what extent has the researcher accounted for false positives and false negatives in the output of the model}. \emph{Dissemination}: report exact model version and whether pretrained \pc{aoir2019}{44}{What exact version of the model did the researcher use, was this model pre-trained}; reproducibility means describing the experiment so that another person could get at least similar results \pc{aoir2019}{44}{The scientific requirements of reproducibility and replicability demand from researchers to describe their experiment in such a way that another person could achieve at least similar results}; transparency means explaining ``the ends, means, and thought processes'' behind code and model \pc{aoir2019}{45}{a richer notion of transparency asks researchers to explain the ends, means, and thought processes that went into developing the code, the model, and how the resulting research data was shaped}, and the NeurIPS reproducibility checklist is cited \pc{aoir2019}{45}{The NeurIPS conference for example requires completion of a reproducibility checklist}. Downstream: what are the consequences of erroneous labelling or categorisation \pc{aoir2019}{46}{What could be the downstream consequence for data subjects for erroneous identifications, labelling, or categorization}. \emph{Close}: destruction or archiving \pc{aoir2019}{47}{Will the data be destroyed at a specific date}.

\cf{Limits (card)} A professional-association guideline, not an ethical approval and not binding; written mainly for research on people and social media (not for analysis of company documents or app stores); question lists, no decision rules or thresholds; some content is thin (the preview says example consent forms are in earlier versions); the AI companion predates general-purpose LLMs (2020 text, ``neural networks'' as the example) so LLM-specific issues such as prompt logging, vendor data retention and non-determinism are not addressed; US and Europe-centred examples although cross-cultural pluralism is stressed.

\cf{Use for this thesis} (1) A recognised frame for the ethics paragraph: the thesis analyses published policies, store listings and the behaviour of apps, so human-subject consent is not the main issue, but the stage-based questions (data minimisation, storage, retention, sharing, destruction) apply to scraped store data, screenshots and any user-review text; keep personal data (usernames, reviews) out of the corpus or pseudonymise it. (2) Terms of service and scraping: the guidelines flag that platforms may prohibit automated collection and that the legal status is unsettled; the thesis should state which store terms apply to its scrapers (\texttt{scripts/appstore-scraper}, \texttt{scripts/googleplay-scraper}), the request rate, and why the collection is limited to what is needed; ask the supervisor whether the university requires a statement. (3) LLM-assisted coding (stage L): record exact model version, prompt, settings and date; keep a human check on outputs and say how disagreements are overridden and documented; report false positives and negatives; do not treat LLM labels as ground truth (these map to the companion's automation and bias questions). (4) ``Ethics is method'': log ethical decisions in the Decision Log as they arise. (5) Reflexivity and downstream harm: state that coding of harms can itself mislabel an app or developer, so wording of findings should avoid asserting wrongdoing that was not observed. (6) The NESH remark links to contextual integrity, which the thesis already uses.
\end{papercard}
